\documentclass[11pt,a4paper]{article}

%% PACHETE MATEMATICE
\usepackage{amsmath,amssymb,amsfonts}
\usepackage{amsthm}
\usepackage{mathpazo}
\usepackage{eulervm}
\usepackage{enumerate}
\usepackage{color}
\usepackage{graphicx}
\usepackage[all]{xy}
\usepackage{xcolor}
\usepackage{framed}
\usepackage[unicode]{hyperref}
\setcounter{MaxMatrixCols}{20}
\usepackage{multicol}
%\usepackage[romanian]{babel}


%% ASPECT
\linespread{1.05}
\usepackage[T1]{fontenc}
\usepackage[utf8x]{inputenc}
\usepackage[margin=2cm]{geometry}
% \usepackage[color]{showkeys}
\usepackage{anyfontsize}
\usepackage{titlesec}
\titleformat*{\section}{\large\bfseries}

\usepackage{fancyhdr}
\pagestyle{fancy}
\rhead{\small \color{gray} Notițe de Adrian Manea}
\lhead{\small \color{gray} ETTI-AG, 2017---2018, A. Niță \& A. Oprina}


\usepackage[autostyle = false, style = german]{csquotes}
\usepackage[romanian]{babel}
\newcommand\qq{\enquote}

%% COMENZILE MELE
\renewcommand*{\proofname}{Dem.:}
\newcommand\bb{\mathbb}
\newcommand\dr{\mathrm}
\newcommand\kal{\mathcal}
\newcommand\fk{\mathfrak}
\newcommand\op{\oplus}
\newcommand\ot{\otimes}
\newcommand\ds{\displaystyle}
\newcommand\id{\indent}
\newcommand\nid{\noindent}
\newcommand\seq{\subseteq}
\newcommand\sneq{\subsetneq}
\newcommand\speq{\supseteq}
\newcommand\spneq{\supsetneq}
\newcommand\rar{\rightarrow}
\newcommand\Rar{\Rightarrow}
\newcommand\xrar{\xrightarrow}
\newcommand\lar{\leftarrow}
\newcommand\Lar{\Leftarrow}
\newcommand\xlar{\xleftarrow}
\newcommand\lrar{\leftrightarrow}
\newcommand\Lrar{\Leftrightarrow}
\newcommand\ideal{\trianglelefteq}
\newcommand\ai{\text{ a.\^{i}. }}
\newcommand\sk{(\textit{Schi\c{t}\u{a}}) }
\newcommand\rhup{\rightharpoonup}
\newcommand\rhdn{\rightharpoondown}
\newcommand\lhup{\leftharpoonup}
\newcommand\lhdn{\leftharpoondown}
\newcommand\vid{\emptyset}
\newcommand\wh{\widehat}
\newcommand\ol{\overline}
\newcommand\NN{\mathbb{N}}
\newcommand\RR{\mathbb{R}}
\newcommand\ZZ{\mathbb{Z}}
\newcommand\QQ{\mathbb{Q}}
\newcommand\CC{\mathbb{C}}

%% STILURILE MELE
\newtheoremstyle{dotless}{}{}{\itshape}{}{\bfseries}{:}{ }{}

\theoremstyle{dotless}
\newtheorem{theorem}{Teorem\u{a}}[section]
\newtheorem{proposition}{Propozi\c{t}ie}[section]
\newtheorem{lemma}{Lem\u{a}}[section]
\newtheorem{corollary}{Corolar}[section]
\newtheorem{exercise}{Exerci\c{t}iu}

\newtheoremstyle{dotlessdr}{}{}{}{}{\bfseries}{:}{ }{}
\theoremstyle{dotlessdr}
\newtheorem{example}{Exemplu}[section]
\newtheorem{definition}{Defini\c{t}ie}[section]
\newtheorem{remark}{Observa\c{t}ie}[section]




\begin{document}

\begin{center}
  {\large\textbf{Seminar 7 \\
    Spații euclidiene\\ Completări și exerciții suplimentare}}
\end{center}

\vspace{1cm}


\section{Distanță}
\label{sec:distanta}

În geometria plană, distanța între vectorii $\vec{v}$ și $\vec{w}$ se calculează prin
$|\vec{v} - \vec{w}|$. De fapt, făcînd un pas și mai în spate, ne putem aminti de distanța
între două puncte de pe axă (vectori într-un spațiu 1-dimensional), calculată prin modulul
diferenței acestora.

Așadar, într-un spațiu euclidian arbitrar, putem generaliza această noțiune astfel:

\begin{definition}\label{def:dist}
  Fie $V$ un spațiu euclidian arbitrar și $x, y \in V$. Se definește \emph{distanța}
  între vectorii $x, y$ prin formula:
  \[
    d(x, y) = || x - y ||,
  \]
  unde în membrul drept apare funcția \emph{normă}, indusă în mod unic de produsul scalar.
\end{definition}

\begin{remark}\label{rk:dist-met}
  Se poate arăta că, dacă gîndim $V$ ca un spațiu topologic, atunci funcția distanță de mai sus
  satisface proprietățile unei \emph{metrici}. Așadar, spațiul euclidian devine chiar spațiu metric.
\end{remark}

%%%%%%%%%%%%%%%%%%%%%%%%%%%%%%%%%%%%%%%%%%%%%%%%%%%%%%%%%%%%%%%%%%%%%%


\section{Morfismul adjunct. Matricea lui Gram}
\label{sec:adj-gram}

Ne referim la următoarea definiție:

\begin{definition}\label{def:adj}
  Fie $V$ un spațiu euclidian și $f : V \to V$ o aplicație liniară.

  Se numește \emph{adjuncta aplicației $f$} aplicația $f^\ast : V \to V$, cu proprietatea:
  \[
    \langle f(x), y \rangle = \langle x, f^\ast(y) \rangle, \forall x,y \in V.
  \]
\end{definition}

Modul în care se calculează adjuncta aplicației folosește \emph{matricea lui Gram}:

\begin{definition}\label{def:gram}
  Fie $V$ un spațiu euclidian $n$-dimensional și $B = \{b_i\}$ o bază a sa.

  Matricea $G \in M_n(\RR)$, definită prin $G = (g_{ij}) = \big( \langle b_i, b_j \rangle \big)$
  se numește \emph{matricea lui Gram} (sau \emph{\qq{gramiană}}).
\end{definition}

De remarcat faptul că \emph{matricea Gram este inversabilă}, în orice bază și indiferent de
produsul scalar, deoarece, din chiar definiția produsului scalar, diagonala principală
a matricei Gram conține produsul pătratelor normelor elementelor din bază.

De exemplu, pentru spațiul euclidian $V = \RR^2$, cu produsul scalar canonic
\[
  \langle x, y \rangle = x_1 y_1 + x_2 y_2, \forall x = (x_1, x_2), y = (y_1, y_2) \in \RR^2,
\]
matricea gramiană este $I_2$, deoarece $\langle e_i, e_j \rangle = \delta_{ij}$, simbolul
lui Kronecker (nul pentru $i \neq j$ și 1 pentru $i = j$).

Folosind matricea gramiană, putem calcula simplu morfismul adjunct. Deoarece orice vector
dintr-un spațiu vectorial arbitrar se scrie în funcție de vectorii bazei și folosind
proprietățile produsului scalar, se poate arăta că în orice spațiu euclidian $V$ avem:
\begin{equation}
  \label{eq:gram}
  \langle v, w \rangle = v^t \cdot G \cdot w, \forall v, w \in V,
\end{equation}
unde $G$ este matricea lui Gram.

Fie $f : V \to V$ o aplicație liniară, căreia vrem să-i determinăm adjuncta. Fie $A$ matricea
aplicației $f$ într-o bază $B = \{b_i\}$ a lui $V$. Atunci, a căuta $f^\ast$ este echivalent
cu a căuta $A^\ast$, iar definiția morfismului adjunct se poate rescrie cu matrice în forma:
\[
  \langle v, A^\ast w \rangle = \langle Av, w \rangle.
\]

Atunci, folosind relația ~\eqref{eq:gram}, avem că:
\[
  v^t GA^\ast w = (Av)^t Gw \Rightarrow (v^t GA^\ast - v^t A^t G) w = 0.
\]
Cum relația are loc pentru $w$ arbitrar, rezultă:
\[
  v^t GA^\ast - v^t A^t G = 0,
\]
iar această relație are loc pentru orice $v$, deci $GA^\ast = A^t G$, de unde, în fine:
\[
  A^\ast = G^{-1} A^t G.
\]


%%%%%%%%%%%%%%%%%%%%%%%%%%%%%%%%%%%%%%%%%%%%%%%%%%%%%%%%%%%%%%%%%%%%%%

\section{Exerciții}
\label{sec:exc}



1. Fie subspațiul:
\[
  V = \{(x_1, x_2, x_3, x_4) \in \RR^4 \mid
  \begin{cases}
    x_1 - 2x_2 + 2x_3 + 2x_4 &= 0 \\
    2x_1 - 4x_2 + x_3 - 2x_4 &= 0 \\
    x_1 - 2x_2 + 3x_3 + 4x_4 &= 0
  \end{cases} \}
\]
al spațiului euclidian $\RR^4$.

Să se determine $V^\perp$ și să se verifice că $V \oplus V^\perp \simeq \RR^4$.

\vspace{1cm}

2. Fie spațiul vectorial $V = \CC^2$ și $B = \{e_1, e_2\}$ baza canonică.
Definim produsul scalar:
\[
  \langle x, y \rangle = 2x_1 \overline{y}_1 - ix_1 \overline{x}_2 + ix_2 \overline{y}_1 + %
  x_2 \overline{y}_2,
\]
unde $x = (x_1, x_2), y = (y_1, y_2)$, iar $\overline{x}$ notează conjugatul complex.

Fie $f : V \to V$ dată de matricea $M_f^B = \begin{pmatrix} -1 & i \\ i & 2 \end{pmatrix}$.

Să se determine endomorfismul adjunct $f^\ast$.

\vspace{1cm}

3. Fie spațiul vectorial $V = \{(x_1, x_2, x_3) \mid 3x_1 + x_2 - x_3 = 0 \}$.
Să se determine $V^\perp$ și coordonatele vectorului $v = (1, -1, 4)$ în $V$ și $V^\perp$.

\vspace{1cm}

4. Fie spațiul vectorial $M_2(\RR)$ și produsul scalar dat de:
\[
  \langle A, B \rangle = \mathrm{tr}(A^t B).
\]

\begin{enumerate}[(a)]
  \item Să se determine matricele de normă 1, ortogonale simultan pe:
\[
  C = \begin{pmatrix}
    1 & 0 \\
    2 & 1
  \end{pmatrix}, \quad
  D = \begin{pmatrix}
    1 & -1 \\
    1 & 0
  \end{pmatrix}
\]
\item Să se completeze mulțimea $\{B, C\}$ la o bază a lui $M_2(\RR)$ și să se
  ortonormeze baza, în raport cu produsul scalar definit mai sus.
\end{enumerate}

\vspace{1cm}

5. În fiecare din spațiile euclidiene $V$ de mai jos, determinați adjunctul morfismului $f$ definit,
relativ la produsul scalar indicat:
\begin{enumerate}[(a)]
\item $V = \RR^2, \langle x, y \rangle = 2x_1 y_1 + 3x_2 y_2, f(x, y) = (2x + y, x - y)$;
\item $V = \RR^2, \langle x, y \rangle = x_1y_1 + 2x_1y_2 + 2x_2 y_1 + 3x_2 y_2$, iar aplicația
  $f(x, y) = (3x, -y)$;
\item $V = \RR_2[X], \langle p, q \rangle = \ds\sum_{k = 0}^2 (k!)^2 a_k b_k$, unde
  $p = a_0 + a_1 X + a_2 X^2$, iar $q = b_0 + b_1 X + b_2 X^2$. Definim aplicația:
  $f(a + bX + cX^2) = (a + b) + (b + c)X + (a + c)X^2$.
\end{enumerate}
\end{document}
